% !PRINTED 650
\weberpage{650}

We can then prove the theorem:

1. \gesperrt{If the group of an Abelian field is decomposable, then the field is also composite.}

To prove this, let us assume that the group $\mathfrak{G}$ of the
field $\Omega$ is decomposable according to the formula (1), and that $m, a, b$
are the degrees of $\mathfrak{G}, \mathfrak{A}, \mathfrak{B}$. We take a number $\xi$ of the field $\Omega$
\gesperrt{belonging} to the group $\mathfrak{B}$, which therefore passes into $a$ different
conjugate values under the substitutions of $\mathfrak{G}$, and likewise
a $b$-valued number $\eta$ belonging to $\mathfrak{A}$. We can
then form a number
\[
x = \alpha \xi + \beta \eta
\]
with \gesperrt{rational} numerical coefficients $\alpha, \beta$, which has $m$ different
conjugate values, and then $\Omega = R(x)$. Hence
$\Omega$ is composed of $R(\xi)$ and $R(\eta)$ (Bd. I, \S.\,143).

If we now recall the representation of an Abelian group by
a basis derived in \S.\,9, 10 of this volume, then repeated
application of what has just been proved yields:

2. \gesperrt{Every Abelian field $\Omega$ can be composed of
fields whose degree is a prime power
and whose group is cyclic. The degrees
of these composing fields are the invariants
of the group of $\Omega$.}

These theorems are supplemented by the following one:

3. \gesperrt{An Abelian field of prime-power degree
with cyclic group is simple.}

If the group $\mathfrak{G}$ is cyclic, then its elements
can be represented by the powers of a basis, $G^\gamma$, when $\gamma$
runs through a complete system of residues with respect to the modulus $m$. We
obtain all divisors $\mathfrak{A}$ of $\mathfrak{G}$ represented by
\[
\mathfrak{A} = G^{b\gamma},
\]
when $b$ is a divisor of $m$ and $\gamma$ runs through a complete system of residues
with respect to the modulus $a = m : b$.

If
\[
\mathfrak{A}' = G^{b'\gamma}
\]
is a second divisor of $\mathfrak{G}$ of degree $a'$, and if
\[
b' \leqq b, \quad a' \geqq a,
\]
then, if we now assume that $m$, and hence also
$a, b, a', b'$, are powers of a prime $q$, $b'$ is a divisor of $b$,
and consequently $\mathfrak{A}$ is a divisor of $\mathfrak{A}'$.
